\section{Limitations and Conclusion}
\label{sec:conclusion}

\paragraph{Limitations.}
\method relies on ground-truth crops on all four benchmarks, as LiePose does on T-LESS; with detected boxes and masks the inputs of
the world model and of the denoiser change, and we have not evaluated this setting. The world model needs a bank of posed training
views of each object and the object's identity at test time; it does not apply to unseen objects. Its precision is bounded by what
frozen, pooled DINOv2 tokens can distinguish (\cref{tab:wm-disc}), and the retrieval and world-model queries at every step make sampling
slower than LiePose (\cref{tab:runtime}). Geometry tokens require the camera intrinsics and the crop box. Finally, our runs use 64 noisy
poses per training image instead of LiePose's 256. \TODO{revise once the results are in; add any failure mode seen in \cref{fig:qual}.}

\paragraph{Conclusion.}
We moved render-and-compare into SE(3) diffusion: a frozen world model predicts, at every denoising step and for every sample, the
DINOv2 tokens of the object at the current hypothesis, and the denoiser conditions on their comparison with the observation. The same
prediction gives an energy that ranks samples without ground truth. Because the diffusion formulation is unchanged and the control
differs only by the predicted tokens, the experiments isolate what the comparison contributes. \TODO{one or two sentences of findings,
written after the tables are filled.}
